\subsection{AXIOMS}

We have already seen that the specific signs determine the terms and the relations in a theory \(\mathfrak{G}\) . To construct \(\mathfrak{G}\) , we proceed as follows:

\begin{enumerate}
\def\labelenumi{(\arabic{enumi})}
\item
  First we write down a certain number of relations in \(\mathfrak{G}\) ; these are called the explicit axioms of \(\mathfrak{G}\) . The letters which appear in the explicit axioms are called the constants of \(\mathfrak{G}\) .
\item
  We lay down one or more rules\footnote{For the sake of brevity, these rules are expressed by using the symbols mentioned in §1, no. 1 (and especially bold-faced italic letters); but it would be easy to avoid the use of these symbols completely in the formulation of the rules (see §3, no. 1, note (*) on p. 28).}, called the schemes of \(\mathfrak{C}\) , which must have the following properties: (a) the application of such a rule \(\mathcal{R}\) furnishes a relation in \(\mathfrak{C}\) ; (b) if \(T\) is a term in \(\mathfrak{C}\) , if \(x\) is a letter, and if \(R\) is a relation in \(\mathfrak{C}\) constructed by applying the scheme \(\mathcal{R}\) , then the relation \((T|x)R\) can also be constructed by applying \(\mathcal{R}\) .
\end{enumerate}

In all the cases we envisage, the verification of these conditions is always easy.

Every relation constructed by applying a scheme of \(\mathfrak{G}\) is called an implicit axiom of \(\mathfrak{G}\) .

Intuitively, the axioms represent either self-evident assertions or else hypotheses from which one wishes to draw consequences. The constants represent well-defined objects for which the properties expressed by the explicit axioms are supposed to be true. On the other hand, if the letter x is not a constant, it represents a completely undetermined object; if a property of the object x is assumed to be true by means of an axiom, then this axiom is necessarily implicit, so that the property remains true for any object T.

